\chapter{Gradings}

\section{Overview}

This chapter records the basic grading data used throughout the project:
the parity group, the class of grading indices, graded modules, and graded
$R$-linear quivers.

\begin{definition}[Parity group]
\lean{AInfinityTheory.Parity}
\label{AInfinityTheory.Parity}
\leanok
The type of parities is $\mathbb{Z}/2\mathbb{Z}$, implemented as \texttt{ZMod 2}.
\end{definition}

\begin{definition}[Grading]
\lean{AInfinityTheory.Grading}
\label{AInfinityTheory.Grading}
\leanok
\uses{AInfinityTheory.Parity}
A grading on a type $\beta$ consists of an additive commutative group structure
on $\beta$, together with homomorphisms
\[
  \phi : \mathbb{Z} \to \beta,
  \qquad
  \sigma : \beta \to \mathbb{Z}/2\mathbb{Z},
\]
such that $\sigma(\phi(n))$ is the parity class of $n$ for every integer $n$.
\end{definition}

\begin{definition}[Integer shift]
\lean{AInfinityTheory.shift_ofInt}
\label{AInfinityTheory.shift_ofInt}
\leanok
\uses{AInfinityTheory.Grading}
For a grading index $\beta$, the definition \lean{AInfinityTheory.shift_ofInt} is the image of
$n \in \mathbb{Z}$ under the distinguished map $\mathbb{Z} \to \beta$.
\end{definition}

\begin{definition}[Graded $R$-module]
\lean{AInfinityTheory.GradedRModule}
\label{AInfinityTheory.GradedRModule}
\leanok
\uses{AInfinityTheory.Grading}
For a commutative ring $R$, a graded $R$-module is a $\beta$-indexed family of
$R$-modules, implemented as a graded object in \texttt{ModuleCat R}.
\end{definition}

\begin{definition}[$R$-linear graded quiver]
\lean{AInfinityTheory.RLinearGQuiver}
\label{AInfinityTheory.RLinearGQuiver}
\leanok
\uses{AInfinityTheory.GradedRModule}
An $R$-linear graded quiver on an object type $\mathrm{Obj}$ assigns to each pair
of objects $(X,Y)$ a graded $R$-module of morphisms from $X$ to $Y$.
\end{definition}
